% Ad Familiares 10.11 --- headnote
% ad-familiares-10-11, late April 43 BC, among the Allobroges

\textit{L. Munatius Plancus to Cicero, written in the territory
of the Allobroges at the end of April 43 BC --- Perseus dateline
\textit{Scr. in Allobrogibus ex. m. Apr. a. 711 (43)}. Plancus
has crossed the Rhone and is on the road to Mutina, his brother
sent ahead with three thousand horse, when news reaches him in
mid-march of the battle at Forum Gallorum and the relief of D.
Brutus from the siege of Mutina. The letter is the report he
sends back from the turning point: Antony is on the run, the
question is now where he is heading, and the answer --- that he
has only Lepidus and Lepidus's army left to him --- is precisely
the answer that lands the whole problem on Plancus's desk.}

\textit{The opening section returns the warmth of Cicero's
\textit{ad Fam.} 10.10, the previous letter in the
correspondence and the one Plancus is now answering. Cicero had
done him \textit{magna officia} ``great services'' in the Senate
--- proposals with unlimited gifts (the first decree of honours,
moved while the case was still hypothetical), later motions
trimmed to the moment, an unbroken speech on his behalf, and
public quarrels with detractors. Plancus disclaims any prospect
of repaying these and accepts Cicero's own formula: it is
enough that he remember them. The closing image of Cicero as
the man who must \textit{tuum munus tuere}, ``look after your
own creation,'' is more candid than it sounds; Plancus is
asking Cicero in the Senate to defend the loyalist Plancus he
has just publicly vouched for.}

\textit{The military middle section is the practical heart of
the letter. Plancus has halted in the country of the Allobroges
to keep his options open: if Antony comes through stripped of
forces, Plancus can handle him alone even should Lepidus's army
take him in; if Antony brings serious forces, and especially if
the Tenth Legion (Caesar's old veterans, whom Plancus had been
instrumental in calling back to the colours) reverts to its
former allegiance, then Plancus needs the Senate's armies sent
across to him from Italy. The third section is the political
overlay: Plancus is working on Lepidus, with his brother and
Laterensis and Furnius as intermediaries, and is willing to
work even with personal enemies for the sake of the cause. The
last clause --- ``and if I gain nothing by it, I shall still
satisfy you, none the less, with the greatest spirit and
perhaps with the greater glory to myself'' --- carries the
shape of a man who already sees how the gambit may fail. Within
weeks Lepidus will join Antony, and Plancus, after a further
hesitation, will follow.}
